\documentclass[11pt,a4paper]{article}

\usepackage[a4paper,margin=1in]{geometry}
\usepackage{booktabs}
\usepackage{enumitem}
\usepackage{hyperref}
\usepackage{xcolor}
\usepackage{listings}
\usepackage{longtable}
\usepackage{titlesec}

\hypersetup{
  colorlinks=true,
  linkcolor=blue!50!black,
  urlcolor=blue!50!black,
  citecolor=blue!50!black
}

\lstdefinestyle{astra}{
  basicstyle=\ttfamily\small,
  breaklines=true,
  frame=single,
  columns=fullflexible,
  showstringspaces=false,
  tabsize=2
}

\titleformat{\section}{\large\bfseries}{\thesection}{0.8em}{}
\titleformat{\subsection}{\normalsize\bfseries}{\thesubsection}{0.8em}{}

\title{\textbf{AstraQPU}\\
\large A Virtual QPU Runtime for Architecture-Aware Scheduling, Control, and Hardware-in-the-Loop Execution}
\author{Project Development Report}
\date{Initial development record: May 2026}

\begin{document}
\maketitle

\begin{abstract}
AstraQPU is a full-stack quantum systems software project focused on the runtime architecture around quantum processor control. The project models a quantum processing unit as a constrained hardware target with native instructions, timing rules, control channels, readout latency, calibration state, execution traces, and hardware-in-the-loop control. This report tracks the project concept, design decisions, implemented components, verification steps, and near-term roadmap.
\end{abstract}

\tableofcontents
\newpage

\section{Project Intent}

The central objective of AstraQPU is to prototype the software architecture around how a quantum processor could be controlled, scheduled, debugged, and virtualized. The project is intentionally positioned as quantum systems infrastructure rather than a classroom circuit simulator.

\subsection{Public Positioning}

\begin{quote}
AstraQPU is a virtual QPU runtime for architecture-aware scheduling, control, and hardware-in-the-loop execution.
\end{quote}

The project message is:

\begin{quote}
I am building the software architecture around how a quantum processor could be controlled, scheduled, debugged, and virtualized.
\end{quote}

\subsection{Design Boundary}

AstraQPU does not begin with state-vector simulation as the main identity. Instead, it begins with the software stack that would exist around a QPU:

\begin{itemize}[leftmargin=*]
  \item instruction ingestion and validation
  \item architecture specification
  \item native operation scheduling
  \item hardware constraints
  \item control-plane communication
  \item timing traces
  \item telemetry and fault reporting
  \item future dashboard and debugging tools
\end{itemize}

\section{Name and Component Language}

The selected project name is \textbf{AstraQPU}. The name combines a Sanskrit-derived root with a direct technical suffix. The internal naming system is:

\begin{longtable}{lll}
\toprule
Name & Role & Meaning in the Project \\
\midrule
AstraQPU & Main project & Virtual QPU runtime and architecture platform \\
Sutra & Instruction stream / IR & A structured thread of executable operations \\
Yantra & HAL / hardware control & Instrumentation and control mechanism \\
Niyantra & Scheduler / runtime control & Regulation, timing, and execution governance \\
Akasha & Virtual backend & Abstract execution space \\
Tejas & Telemetry engine & Signals, status, and device observability \\
\bottomrule
\end{longtable}

\section{Architecture Overview}

The first implementation follows a timing-first runtime pipeline:

\begin{lstlisting}[style=astra]
Sutra instruction stream
  -> AstraQPU IR
  -> architecture validation
  -> constraint-aware scheduler
  -> backend execution
  -> execution trace
\end{lstlisting}

\subsection{Runtime Layers}

\begin{enumerate}[leftmargin=*]
  \item \textbf{Frontend}: parses a low-level Sutra instruction stream.
  \item \textbf{IR}: represents operations, qubits, classical bits, waits, barriers, and register updates.
  \item \textbf{Architecture model}: defines qubits, topology, native gates, durations, channels, readout latency, and calibration data.
  \item \textbf{Scheduler}: assigns start times under qubit, channel, topology, and latency constraints.
  \item \textbf{Backend}: executes scheduled streams through either a virtual backend or a serial MCU backend.
  \item \textbf{Trace layer}: records machine-readable execution events.
\end{enumerate}

\section{Repository Structure}

The initial repository layout is:

\begin{lstlisting}[style=astra]
astraqpu/
  astraqpu/
    arch/
    backends/
    frontend/
    ir/
    protocol/
    runtime/
    scheduler/
    cli.py
  docs/
    architecture.md
    instruction-set.md
    main.tex
    qpu-architecture-spec.md
    roadmap.md
    serial-protocol.md
  examples/
    arch/
    bell.aqis
    channel_overlap.aqis
  firmware/
    arduino/
    esp32/
  tests/
\end{lstlisting}

\section{Implemented Development Steps}

\subsection{Step 1: Project Foundation}

The first development pass created a local Git repository under the project workspace and established a professional baseline:

\begin{itemize}[leftmargin=*]
  \item README with project positioning
  \item MIT license
  \item Python package metadata
  \item line-ending normalization through \texttt{.gitattributes}
  \item clean Git baseline on the \texttt{main} branch
\end{itemize}

\subsection{Step 2: Architecture Specification}

A tiny 3-qubit superconducting-style mock target was added at:

\begin{lstlisting}[style=astra]
examples/arch/tiny_3q.json
examples/arch/tiny_3q.yaml
\end{lstlisting}

The architecture specifies:

\begin{itemize}[leftmargin=*]
  \item clock tick size
  \item qubit coherence metadata
  \item physical couplings
  \item native gates
  \item gate durations
  \item readout duration and latency
  \item control channel templates
  \item calibration values
\end{itemize}

\subsection{Step 3: Sutra Instruction Stream}

The first instruction format supports:

\begin{lstlisting}[style=astra]
DECLARE_QUBITS n
DECLARE_BITS n
PREP q
GATE name q...
MEASURE q -> c
WAIT duration
BARRIER q...
SET_REG name value
END
\end{lstlisting}

An example Bell program was added:

\begin{lstlisting}[style=astra]
DECLARE_QUBITS 2
DECLARE_BITS 2

PREP q0
PREP q1
GATE H q0
GATE CX q0 q1
MEASURE q0 -> c0
MEASURE q1 -> c1
END
\end{lstlisting}

\subsection{Step 4: Constraint-Aware Scheduler}

The scheduler currently enforces:

\begin{itemize}[leftmargin=*]
  \item declared qubits must exist in the architecture
  \item native operations must exist in the architecture
  \item two-qubit gates must follow the coupling topology
  \item qubits cannot be reused before their previous operation and latency window ends
  \item channels cannot be reused before their operation duration ends
  \item explicit waits and barriers affect scheduling
  \item operation durations are rounded to the architecture clock tick
\end{itemize}

\subsection{Step 5: Akasha Virtual Backend}

The virtual backend converts a scheduled program into an execution trace with events such as:

\begin{itemize}[leftmargin=*]
  \item \texttt{instruction\_start}
  \item \texttt{instruction\_end}
  \item \texttt{latency\_complete}
\end{itemize}

This provides the first runtime proof: AstraQPU can compile a low-level QPU instruction stream into a hardware-constrained timing trace.

\subsection{Step 6: Yantra Serial Protocol}

A JSON Lines serial protocol was introduced for hardware-in-the-loop control. The host emits:

\begin{lstlisting}[style=astra]
{"type":"HELLO","proto":"astraqpu.serial.v0","host":"astraqpu"}
{"type":"LOAD","proto":"astraqpu.serial.v0","job_id":"bell_001","architecture":"tiny_3q_superconducting_mock","tick_ns":10,"instruction_count":6}
{"type":"INST","proto":"astraqpu.serial.v0","job_id":"bell_001","id":1,"t_ns":0,"op":"prep","qubits":["q0"],"duration_ns":20,"latency_ns":0,"channels":["drive:q0"]}
{"type":"RUN","proto":"astraqpu.serial.v0","job_id":"bell_001"}
\end{lstlisting}

The MCU is expected to return:

\begin{lstlisting}[style=astra]
{"type":"HELLO","proto":"astraqpu.serial.v0","device":"arduino","tick_ns":1000}
{"type":"ACK","job_id":"bell_001","command":"LOAD"}
{"type":"EVT","job_id":"bell_001","event":"instruction_start","id":1,"op":"prep","t_ns":0}
{"type":"RESULT","job_id":"bell_001","bits":{}}
{"type":"DONE","job_id":"bell_001"}
\end{lstlisting}

\subsection{Step 7: Arduino Control Unit}

The Arduino firmware now implements a dependency-free protocol emulator. It accepts:

\begin{itemize}[leftmargin=*]
  \item \texttt{HELLO}
  \item \texttt{LOAD}
  \item \texttt{INST}
  \item \texttt{RUN}
\end{itemize}

It stores a fixed-size instruction buffer, emits ACK/NACK responses, emits timing events, and returns RESULT and DONE messages.

\section{Command-Line Interface}

The current CLI entry point is:

\begin{lstlisting}[style=astra]
python -m astraqpu.cli
\end{lstlisting}

Available workflows:

\begin{lstlisting}[style=astra]
python -m astraqpu.cli validate examples/arch/tiny_3q.json
python -m astraqpu.cli compile examples/bell.aqis --arch examples/arch/tiny_3q.json --out build/bell.scheduled.json
python -m astraqpu.cli run examples/bell.aqis --arch examples/arch/tiny_3q.json --backend virtual
python -m astraqpu.cli serial-dump examples/bell.aqis --arch examples/arch/tiny_3q.json --job-id bell_001
python -m astraqpu.cli run examples/bell.aqis --arch examples/arch/tiny_3q.json --backend serial --port COM5
\end{lstlisting}

\section{Verification}

The current repository includes standard-library unit tests. No external test dependency is required.

\begin{lstlisting}[style=astra]
python -m unittest discover -s tests
\end{lstlisting}

The initial verified behaviors are:

\begin{itemize}[leftmargin=*]
  \item Sutra parsing
  \item duration parsing
  \item Bell program scheduling
  \item readout latency handling
  \item topology rejection for uncoupled two-qubit operations
  \item virtual trace emission
  \item serial JSON Lines generation and parsing
\end{itemize}

\section{Current Engineering Status}

\begin{longtable}{lll}
\toprule
Area & Status & Notes \\
\midrule
Architecture spec & Implemented & JSON supported, YAML optional with PyYAML \\
Sutra parser & Implemented & Low-level v0 instruction stream \\
Scheduler & Implemented & Timing, topology, qubits, channels, latency \\
Virtual backend & Implemented & Produces trace JSON \\
Serial protocol & Implemented & JSON Lines host protocol and parser \\
Serial backend & Implemented & Requires pyserial and a board port \\
Arduino firmware & Initial implementation & Dependency-free protocol emulator \\
ESP32 firmware & Planned & Recommended as primary v0.2 target \\
Dashboard & Planned & Timeline and telemetry viewer later \\
\bottomrule
\end{longtable}

\section{Next Milestones}

\subsection{v0.2 Hardware-in-the-Loop Completion}

\begin{itemize}[leftmargin=*]
  \item test Arduino firmware on a physical board
  \item confirm host serial run receives DONE and RESULT
  \item add calibration register commands
  \item add fault flag messages
  \item improve MCU timing model on ESP32
\end{itemize}

\subsection{v0.3 Observability}

\begin{itemize}[leftmargin=*]
  \item generate channel occupancy timelines
  \item generate simple HTML or web dashboard
  \item compare scheduled trace against MCU-reported trace
  \item record telemetry streams
\end{itemize}

\subsection{v0.4 Compiler Frontend}

\begin{itemize}[leftmargin=*]
  \item parse an OpenQASM subset
  \item lower circuit-level gates to native gates
  \item add mapping passes
  \item add calibration-aware scheduling
\end{itemize}

\section{Conclusion}

AstraQPU has been initialized as a serious quantum systems infrastructure project. The current codebase already demonstrates the core identity of the project: a QPU program can be parsed, checked against a hardware-inspired architecture, scheduled under timing and topology constraints, executed virtually, converted into a serial control stream, and prepared for physical microcontroller-backed execution.

\end{document}

